\section{Arithmetic separation and counting}
\label{sec:C-density}

The period criteria above concern a single algebraic fiber.  We now use
the variation of that fiber to obtain unconditional restrictions on the
set of possible identities.  Two different principles apply.  Near the
cusp, independent Gauss systems and the Bombieri--Andr\'e specialization
theorem separate identities whose multipliers have bounded degree.
On a compact parameter interval, symmetric-square monodromy and
Habegger--Pila counting give a subpower bound when all arithmetic heights
are bounded.  Neither argument assumes the period hypothesis of
Theorem~\ref{thm:main-conditional}.

For $d\geq1$, let $\mathscr E_d$ be the set of triples
$(\ell,\epsilon,t)$, with $1\leq\ell\leq4$, $\epsilon\in\{\pm1\}$,
and $t\in\mathbb N$, for which there is a primitive standard identity
\begin{equation}\label{rev:eq:exceptional-occurrence}
 (A+B\thetaop)F_{s_\ell}\!\left(\frac{\epsilon R_\ell}{t}\right)
       =\frac{\alpha}{\pi},
 \qquad [\mathbb Q(\alpha^2):\mathbb Q]\leq d.
\end{equation}
Thus $t$ is the absolute denominator, while $(\ell,\epsilon)$ retains
the level and its sign.  By Theorem~\ref{thm:unconditional}, each such
triple determines at most one primitive coefficient pair up to sign.

\begin{theorem}[Separation in absolute denominator]
\label{thm:C-absolute-gaps}
Fix $d\geq1$.  For every integer $H\geq0$, there exists $T_H$ such that
\[
 \#\{(\ell,\epsilon,t)\in\mathscr E_d:
                      M\leq t\leq M+H\}\leq1
 \qquad(M\geq T_H).
\]
Equivalently, for fixed levels $\ell_1,\ell_2$, signs
$\epsilon_1,\epsilon_2$, and integer $D\geq0$, there are only finitely many
$t$ for which
\[
 (\ell_1,\epsilon_1,t),\qquad
 (\ell_2,\epsilon_2,t+D)\in\mathscr E_d,
\]
unless $D=0$ and $(\ell_1,\epsilon_1)=(\ell_2,\epsilon_2)$.
In particular, distinct absolute denominators have gaps tending to
infinity, jointly across all four levels and both signs.
\end{theorem}

The theorem imposes no height bound on $A,B$, or $\alpha$, and no fixed
field of definition for $\alpha$.  The square in its degree hypothesis
is natural: taking the squared ratio of two identities eliminates
their common transcendental factor.

\subsection{The Gauss systems and their products}

Write
\[
 \varphi_s(z)={}_2F_1\!\left(\frac{s}{2},\frac{1-s}{2};1;z\right),
 \qquad F_s=\varphi_s^2,\qquad G_s=\thetaop F_s,
\]
and let $\mathcal V_s$ denote the rank-two local system of solutions of
\begin{equation}\label{eq:C-gauss-operator}
 z(1-z)y''+\left(1-\frac32z\right)y'
                      -\frac{s(1-s)}4y=0.
\end{equation}
Clausen's identity identifies the third-order system of $F_s$ with
$\operatorname{Sym}^2\mathcal V_s$.  Its arithmetic and geometric
roles will be kept distinct: the Taylor coefficients make $F_s,G_s$
$G$-functions, whereas monodromy controls their functional relations.

\begin{lemma}[Product monodromy]\label{rev:lem:product-monodromy}
The connected algebraic monodromy of $\mathcal V_s$ is
$\operatorname{SL}_2$ for every standard signature $s$.
Let
\[
 r_i=\epsilon_iR_{\ell_i},\qquad s_i=s_{\ell_i},\qquad
 \rho_1(x)=\frac{r_1}{x},\quad
 \rho_2(x)=\frac{r_2}{x+D},\qquad D\in\mathbb Z_{\geq0}.
\]
Unless $D=0$ and $(\ell_1,\epsilon_1)=(\ell_2,\epsilon_2)$,
the connected algebraic monodromy of
$\rho_1^*\mathcal V_{s_1}\oplus\rho_2^*\mathcal V_{s_2}$ is
\[
                 \operatorname{SL}_2\times\operatorname{SL}_2.
\]
Consequently the four functions
\begin{equation}\label{rev:eq:independent-four}
 F_{s_1}(r_1u),\quad G_{s_1}(r_1u),\quad
 F_{s_2}\!\left(\frac{r_2u}{1+Du}\right),\quad
 G_{s_2}\!\left(\frac{r_2u}{1+Du}\right)
\end{equation}
are algebraically independent over $\mathbb C(u)$.
\end{lemma}

\begin{proof}
Put $a=s/2$ and $b=(1-s)/2$.  Abel's identity gives a Wronskian
proportional to $1/(z\sqrt{1-z})$, so the monodromy determinant has
finite image.  At zero the nontrivial unipotent local monodromy fixes
the line of $\varphi_s$.  At one, the connection formula reads
\[
 \varphi_s(z)=
 \frac{\Gamma(1/2)}{\Gamma(1-a)\Gamma(1-b)}u(z)
 +\frac{\Gamma(-1/2)}{\Gamma(a)\Gamma(b)}
             (1-z)^{1/2}v(z),
\]
where $u(1)=v(1)=1$.  Both coefficients are nonzero.  Hence the
reflection at one moves the fixed line of the unipotent at zero.
The two conjugate one-parameter unipotent groups generate
$\operatorname{SL}_2$; the finite determinant image shows that this
is precisely the identity component.

Let $H$ be the joint algebraic monodromy of the two pullbacks.  Its
identity component projects onto both copies of $\operatorname{SL}_2$.
By the Lie-algebra form of Goursat's lemma, a proper connected subgroup
with these projections would have Lie algebra
\begin{equation}\label{rev:eq:goursat-graph}
 \mathfrak h=\{(X,\phi X):X\in\mathfrak{sl}_2\}
\end{equation}
for an isomorphism $\phi:\mathfrak{sl}_2\to\mathfrak{sl}_2$.
Normality of $H^0$ in $H$ would then require, for every $(M,N)\in H$,
\begin{equation}\label{rev:eq:adjoint-intertwining}
       \phi\operatorname{Ad}(M)=\operatorname{Ad}(N)\phi.
\end{equation}
We rule this out by examining the singularities, including the only
case in which the two sets of singularities coincide.

The finite singular sets are
\[
                  \{0,r_1\},\qquad\{-D,r_2-D\}.
\]
Every point of either set has noncentral local monodromy.  At Gauss
argument one this is the reflection with eigenvalues $1,-1$.  At
infinity the exponent difference is $1/2-s$.  It is nonintegral
unless $s=1/2$.  In that remaining case, scalar monodromy at infinity
would, by the product relation between the three local monodromies,
make the reflection at one a scalar multiple of the unipotent at
zero, which is impossible.  If the finite singular sets differ, a
loop around a point belonging to only one gives $(M,1)$ or $(1,N)$
with a noncentral entry.  Equation~\eqref{rev:eq:adjoint-intertwining}
would make that entry central, a contradiction.

If the sets coincide and the systems are not identical, their common
configuration is necessarily
\[
 \ell_1=\ell_2=\ell,\qquad
 \epsilon_1=-1,\quad\epsilon_2=1,\qquad D=R_\ell.
\]
Indeed $D>0$ then forces $r_1=-D$ and $r_2=D$, and the four radii
are distinct.  A loop about $x=0$ corresponds to Gauss infinity in
the first system and Gauss argument one in the second.  These local
maps are unramified.  For its monodromy pair $(M,N)$, the adjoint
traces are
\begin{equation}\label{rev:eq:collision-traces}
 \begin{split}
 \operatorname{tr}\operatorname{Ad}(N)&=-1,\\
 \operatorname{tr}\operatorname{Ad}(M)
  &=1-2\cos(2\pi s)
    =\begin{cases}
       0,&s=1/6,\\
       1,&s=1/4,\\
       2,&s=1/3,\\
       3,&s=1/2.
      \end{cases}
 \end{split}
\end{equation}
The formula uses only the eigenvalues, so also covers the
nondiagonalizable monodromy when $s=1/2$.  The unequal traces
contradict~\eqref{rev:eq:adjoint-intertwining}.  Thus $H^0$ is the
full product in every permitted case.

To pass from local systems to the selected functions, evaluate at an
ordinary point.  Each group factor moves the nonzero solution column
$(\varphi_s,\varphi_s')$ through a Zariski-dense subset of $\mathbb A^2$.  The
quadratic map
\begin{equation}\label{rev:eq:symmetric-square-projection}
          (v,w)\longmapsto(v^2,2zvw)
\end{equation}
is dominant for $z\ne0$ and sends this column to $(F_s,G_s)$.
The two group factors act independently.  A functional polynomial
relation among the four functions would therefore vanish on
$\mathbb A^4$ at every ordinary point; all its coefficients would
vanish.  Substituting $u=1/x$ proves the asserted independence.
\end{proof}

The same argument has a useful one-system consequence.

\begin{lemma}[Affine functional independence]
\label{lem:C-functional-affine-independence}
For a standard signature, $1,F_s,G_s$ are linearly independent over
$\mathbb C(z)^{\mathrm{alg}}$, on every continued branch.
\end{lemma}
\begin{proof}
Make the coefficients of a proposed relation single-valued on a finite
algebraic cover.  The connected monodromy remains
$\operatorname{SL}_2$, since passing to a finite-index subgroup does
not change its connected Zariski closure.  By the dominant map
\eqref{rev:eq:symmetric-square-projection}, continuation of the
selected solution makes the proposed affine equation vanish at every
point of $\mathbb A^2$.  Its three coefficients are therefore zero.
This is the affine form of irreducibility of
$\operatorname{Sym}^2\mathcal V_s$: its three-dimensional monodromy
orbit cannot satisfy an inhomogeneous first-order equation over an
algebraic extension of the base function field.
\end{proof}

\subsection{Specialization and the norm of a period relation}

We use the following form of the Bombieri--Andr\'e theorem, as stated
in \cite[Theorem~3.13 and Remark~3.14]{BaldiBinyaminiUrbanik2026}.
Its coefficient field is fixed; its polynomial coefficients have no
height restriction.  Both features are essential here.

Recall that a $G$-function over $\mathbb Q$ is a power series satisfying
a linear differential equation over $\mathbb Q(u)$, whose coefficient
sizes and common denominators up to order $n$ grow at most exponentially
in $n$.  The standard $F_s$ have this property: their coefficients become
integers after multiplication by $R_\ell^n$, and their defining
hypergeometric equations are linear.  Differentiation and rational
substitution vanishing at the origin preserve this class.

\begin{theorem}[Bombieri--Andr\'e, archimedean form]
\label{rev:thm:bombieri-andre}
Let $g_1,\ldots,g_m\in\mathbb Q[[u]]$ be $G$-functions, and let
$\mathcal B\subset\mathbb P^1\times\mathbb A^m$ be the Zariski closure
over $\mathbb Q$ of their graph.  Write
\[
 \mu'=\operatorname{trdeg}_{\mathbb Q(u)}
       \mathbb Q(u)(g_i^{(j)}:1\leq i\leq m,\ j\geq0)>0.
\]
There are constants $c_1\geq e$ and $c_2>0$, depending only on the functions, with
the following property.  Let $\delta\geq1$ and let
$\xi\in\mathbb Q\setminus\{0\}$ lie strictly inside all convergence
discs.  Put $h=h(\xi)$, the absolute logarithmic Weil height.  If
\begin{equation}\label{rev:eq:bombieri-andre-bounds}
 \begin{split}
 h&\geq c_1\delta^{\mu'}(1+\log\delta),\\
 |\xi|&\leq
 \exp\!\left[-c_2\delta^{\mu'/(\mu'+1)}
       h^{1/(\mu'+1)}\log\bigl(h^{1/(\mu'+1)}\bigr)\right],
 \end{split}
\end{equation}
there is no polynomial $P\in\mathbb Q[T_1,\ldots,T_m]$ of degree
at most $\delta$ such that
\[
 P(g_1(\xi),\ldots,g_m(\xi))=0,
 \qquad
 \dim(\mathcal B_\xi\cap V(P))<\dim\mathcal B-1.
\]
\end{theorem}

The dimension inequality is the precise meaning of a strongly
nontrivial specialized relation.  The cited theorem is stated for an
arbitrary number field and place; the form above takes
$K=\mathbb Q$ and the real place.  Its statement for degree $\delta$
also gives degree at most $\delta$, by applying it to each smaller
degree.  Nonhomogeneous relations are allowed by adjoining the
constant $1$ to the tuple, as explained in the cited remark.

\begin{lemma}[Reciprocal-integer specialization]
\label{lem:C-reciprocal-specialization}
Let $g_1,\ldots,g_m$ be $G$-functions over $\mathbb Q$, holomorphic at
zero and algebraically independent over $\mathbb Q(u)$.  For each
fixed $\delta\geq1$, their values at $u=1/t$ satisfy no nonzero
polynomial relation over $\mathbb Q$ of total degree at most $\delta$
for all sufficiently large positive integers $t$.
\end{lemma}
\begin{proof}
Functional independence gives
$\mathcal B=\mathbb P^1\times\mathbb A^m$, so every nonzero
specialized relation is strongly nontrivial.  The differential field
has finite transcendence degree $\mu'$: each $G$-function satisfies a
linear differential equation over $\mathbb Q(u)$, so finitely many
derivatives generate it.  At $\xi=1/t$ one has $h=\log t$ and
$|\xi|=e^{-h}$.  For fixed $\delta$, the first inequality in
\eqref{rev:eq:bombieri-andre-bounds} eventually holds; the negative
logarithm required by the second is
$O(h^{1/(\mu'+1)}\log h)=o(h)$.  The theorem applies, uniformly in
the coefficients of the proposed relation.
\end{proof}

\begin{proof}[Proof of Theorem~\ref{thm:C-absolute-gaps}]
Fix the two signed levels and the gap $D$, excluding the identical
case.  Let $\mathbf g=(g_1,g_2,g_3,g_4)$ be the tuple in
\eqref{rev:eq:independent-four}.  The rational hypergeometric
parameters make $F_s$ a $G$-function.  Differentiation and rational
substitution vanishing at zero preserve $G$-functions, hence every
$g_i$ is a $G$-function over $\mathbb Q$, holomorphic at zero.
Lemma~\ref{rev:lem:product-monodromy} gives their algebraic
independence.

Suppose the two identities occur at absolute denominators $t$ and
$t+D$.  Their multipliers are nonzero by
Theorem~\ref{thm:B-value-independent}.  Set
\[
 \beta=\frac{\alpha_1^2}{\alpha_2^2},\qquad
 K=\mathbb Q(\beta),\qquad e=[K:\mathbb Q]\leq d^2,
\]
and regard
\[
 U_1=A_1T_1+B_1T_2,\qquad U_2=A_2T_3+B_2T_4
\]
as linear forms on $\mathbb A^4$.  Cancelling $\pi$ between the
squared identities gives
\[
             (U_1^2-\beta U_2^2)(\mathbf g(1/t))=0.
\]
The descent to the fixed coefficient field is the polynomial norm
\begin{equation}\label{rev:eq:norm-descent}
 Q=N_{K/\mathbb Q}(U_1^2-\beta U_2^2)
   =\prod_{\sigma:K\hookrightarrow\mathbb C}
               (U_1^2-\sigma(\beta)U_2^2)
       \in\mathbb Q[T_1,T_2,T_3,T_4].
\end{equation}
Each factor is nonzero, since $U_1,U_2$ are nonzero forms in disjoint
pairs of variables.  Thus $Q\ne0$, it is homogeneous of degree
$2e\leq2d^2$, and $Q(\mathbf g(1/t))=0$.
Lemma~\ref{lem:C-reciprocal-specialization} excludes this for all
sufficiently large $t$.

For fixed $H$, there are only finitely many ordered pairs of signed
levels and gaps $0\leq D\leq H$.  Taking the largest of their
thresholds proves the interval assertion.  The already established
uniqueness of the primitive coefficient direction identifies repeated
representations of the same signed-level occurrence.
\end{proof}

Notice that the norm in~\eqref{rev:eq:norm-descent} conjugates only
the algebraic coefficient $\beta$.  It does not assert that a Galois
conjugate of a multiplier occurs on the chosen analytic branch.  This
distinction is what allows an unrestricted coefficient field for
$\alpha$ while applying a specialization theorem over $\mathbb Q$.

\subsection{Counting in complete height windows}
\label{sec:C-semirational-counting}

For a real number $y$ and $d\geq1$, use the polynomial coefficient
height
\[
 H_d(y)=\min\{\|q\|_\infty:
     0\ne q\in\mathbb Z[T]\text{ primitive},\quad
     \deg q\leq d,\quad q(y)=0\},
\]
where $\min\varnothing=\infty$ and $\|q\|_\infty$ is the largest
absolute coefficient.  On tuples, $H_d$ is the maximum coordinate
height.  This is the $d$-height used by Habegger--Pila; in particular
$H_d(y)<\infty$ is equivalent to $[\mathbb Q(y):\mathbb Q]\leq d$.

\begin{theorem}[Habegger--Pila, semirational counting]
\label{thm:C-HP-counting}
Let $Z\subseteq\mathbb R^m\times\mathbb R^n$ be definable in a fixed
o-minimal structure, with coordinate projections $\pi_1,\pi_2$.
For every $d\geq1$ and $\varepsilon>0$, there is $c>0$ such that,
whenever $T\geq1$ and $\Sigma\subseteq Z$ satisfy
\[
 H_d(\pi_1P)\leq T\quad(P\in\Sigma),\qquad
                  \#\pi_2(\Sigma)>cT^\varepsilon,
\]
there exists a continuous definable path $\gamma:[0,1]\to Z$ such
that $\pi_1\gamma$ is semialgebraic and real analytic on $(0,1)$,
and $\pi_2\gamma$ is nonconstant.
\end{theorem}

This is \cite[Corollary~7.2]{HabeggerPilaCounting} with a fixed member
of its definable family.  Only the first projection is required to be
semialgebraic; no degree or height restriction is placed on the second.
Our choice of the two projections will use this distinction explicitly.

Fix a compact interval $I\subset[-1,1)$ and define the normalized
incidence surface
\begin{equation}\label{eq:C-normalized-surface}
 X_s=\{(z,a,b)\in(I\setminus\{0\})\times\mathbb R^2:
                      \pi(aF_s(z)+bG_s(z))=1\}.
\end{equation}
The chosen branch is analytic on a neighborhood of $I$, so $X_s$ is
definable in $\mathbb R_{\mathrm{an}}$.  Its nonempty fibers over $z$
are affine lines.  These lines fill the surface and therefore prevent a direct
application of counting only the usual transcendental part.  The
relevant assertion concerns paths that move in the parameter.

\begin{lemma}\label{lem:C-vertical-arcs}
Every continuous semialgebraic path in $X_s$ has constant $z$ coordinate.
\end{lemma}
\begin{proof}
If the parameter varies, semialgebraic cell decomposition and a local
inverse give a nonempty interval on which the path is
$(z,a(z),b(z))$, with $a,b$ real analytic and algebraic over
$\mathbb R(z)$.  Its incidence equation is an affine functional
relation between $1,F_s,G_s$, contrary to
Lemma~\ref{lem:C-functional-affine-independence}.  That lemma allows
arbitrary complex constants; no algebraicity of the constants defining
the semialgebraic path is needed.
\end{proof}

\begin{theorem}[Counting normalized identities]
\label{thm:C-normalized-counting}
For a standard signature $s$, a compact interval $I\subset[-1,1)$,
$d\geq1$, and every $\varepsilon>0$,
\[
 \#\{(z,a,b)\in X_s:H_d(z,a,b)\leq T\}
                    =O_{s,I,d,\varepsilon}(T^\varepsilon).
\]
\end{theorem}
\begin{proof}
Duplicate the parameter by setting
\[
 Z_s=\{((z,a,b),t):(z,a,b)\in X_s,\ t=z\}.
\]
The two projections fit in the commutative square
\[
 \begin{tikzcd}[column sep=large,row sep=large]
 Z_s \arrow[r,"\pi_1"] \arrow[d,"\pi_2"'] &
 X_s \arrow[d,"p"] \\
 I\setminus\{0\} \arrow[r,equal] & I\setminus\{0\},
 \end{tikzcd}
 \qquad p(z,a,b)=z.
\]
If more than $cT^\varepsilon$ distinct parameters occur, the
Habegger--Pila theorem produces a path in $Z_s$ whose first projection
is a semialgebraic path in $X_s$ and whose parameter is nonconstant.
This contradicts Lemma~\ref{lem:C-vertical-arcs}.  Thus the asserted
bound holds for the number of parameters.

Above each counted parameter there is at most one algebraic pair
$(a,b)$.  Indeed, two pairs would give a nonzero algebraic linear
relation between $F_s(z)$ and $G_s(z)$, excluded by
Theorem~\ref{thm:B-value-independent}.  Counting parameters therefore
counts the complete normalized triples.
\end{proof}

\begin{theorem}[Counting complete primitive tuples]
\label{thm:C-primitive-counting}
Fix a standard pair $(s,R)$, $d\geq1$, and $\varepsilon>0$.  The number
of ordinarily convergent primitive identities
\[
 \pi(A+B\thetaop)F_s(R/C)=\alpha,\qquad
 A,B,C\in\mathbb Z,\quad C\ne0,\quad\gcd(A,B)=1,
\]
for which
\[
                  \max\{|A|,|B|,|C|,H_d(\alpha)\}\leq T
\]
is $O_{s,R,d,\varepsilon}(T^\varepsilon)$.  Both simultaneous signs
are included.
\end{theorem}
\begin{proof}
Theorem~\ref{thm:unconditional} excludes $z=1$.  All other ordinarily
convergent integer denominators lie in the fixed interval
\[
                         I_R=[-1,R/(R+1)].
\]
The multiplier is real and nonzero.  Normalize the identity by
\[
                 (z,a,b)=\left(\frac RC,\frac A\alpha,
                                                \frac B\alpha\right)
                 \in X_s.
\]
If $q$ is a primitive polynomial of degree at most $d$ and height at
most $T$ annihilating $\alpha$, then, for $A\ne0$,
$X^d q(A/X)$ is a nonzero integer polynomial annihilating $A/\alpha$.
Removing its content yields
\[
 H_d(A/\alpha)\leq T^{d+1},\qquad
 H_d(B/\alpha)\leq T^{d+1},\qquad
 H_d(R/C)\leq\max(R,T).
\]
The bound for a zero numerator follows from $H_d(0)=1$.  Apply
Theorem~\ref{thm:C-normalized-counting} with exponent
$\varepsilon/(d+1)$.  It gives $O(T^\varepsilon)$ normalized triples.

A normalized triple coming from an integer identity determines a
rational coefficient direction.  Such a direction contains exactly
two primitive integer vectors, and fixing either vector determines $\alpha$.
Thus normalization has fibers of size at most two on the tuples being
counted, proving the result.
\end{proof}

These theorems supply two different forms of arithmetic scarcity.
The separation theorem uses only the degree of $\alpha^2$; the
counting theorem controls the complete tuple in degree and height.
They leave open the possibility of an isolated non-CM identity or an
infinite sequence whose denominators and heights grow sufficiently
rapidly.  In particular, neither conclusion supplies the missing
pointwise CM implication.
